\documentclass[11pt,letterpaper,spanish]{article}

\usepackage{caption}

\usepackage{iftex}
\ifPDFTeX
  % --- Setup for pdfLaTeX ---
  \usepackage[utf8]{inputenc}
  \usepackage[T1]{fontenc}
\else
  % --- Setup for XeLaTeX & LuaLaTeX ---
  \usepackage{fontspec}
  \setmainfont{Latin Modern Roman}
\fi


\usepackage[spanish]{babel}

\usepackage{url}
\usepackage{hyperref}
\usepackage{geometry}
\usepackage{svg}
\usepackage{fancyhdr}
\usepackage{xcolor}
\usepackage{multicol}
\usepackage{enumitem}

%% Paquetes caseros
\usepackage{customcode}

% ==========================================
% MARCADO DE PROCEDENCIA DE LAS JUSTIFICACIONES
% (permite distinguir lo tomado de Explicación.pdf
%  de las propuestas de mejora para poder depurarlas)
% ==========================================
\definecolor{propcolor}{RGB}{0,90,160}

\tcbuselibrary{breakable}

\newcommand{\marcaOriginal}{\textnormal{\small\textcolor{black!55}{(de Explicación.pdf)}}}
\newcommand{\marcaPropuesta}{\textnormal{\small\textcolor{propcolor}{(propuesta)}}}

\newtcolorbox{cajapropuesta}{
	colback=propcolor!4,
	colframe=propcolor,
	boxrule=0.5pt,
	arc=2pt,
	left=5pt, right=5pt, top=5pt, bottom=5pt,
	breakable
}



% Required by biblatex-apa for context-sensitive quotes
\usepackage{csquotes}

% Load biblatex with official APA 7th edition style
\usepackage[style=apa, backend=biber]{biblatex}
\addbibresource{referencias.bib}


\geometry{margin=2cm, top=2.5cm, bottom=2.5cm}

\definecolor{headergray}{gray}{0.3}
\definecolor{rulecolor}{gray}{0.6}

\pagestyle{fancy}
\fancyhf{}
\renewcommand{\headrulewidth}{0.5pt}
\renewcommand{\footrulewidth}{0.5pt}
\fancyhead[L]{\small\textcolor{headergray}{UNIVERSIDAD NACIONAL AUTÓNOMA DE MÉXICO}}
\fancyhead[R]{\small\textcolor{headergray}{\documentoTitulo}}
\fancyfoot[L]{\small\textcolor{headergray}{FACULTAD DE CIENCIAS, 2027-I}}
\fancyfoot[C]{\small\textcolor{headergray}{\thepage}}
\fancyfoot[R]{\small\textcolor{headergray}{FUNDAMENTOS DE BASES DE DATOS}}

% ==========================================
% CONFIGURACIÓN DE VARIABLES (EDITAR AQUÍ)
% ==========================================
\newcommand{\documentoTitulo}{Practica 03}
\newcommand{\documentoSubtitulo}{Modelo Entidad-Relación Extendido}

\newcommand{\profesorNombre}{Gerardo Avilés Rosas}

\newcommand{\ayudanteTeoriaUno}{Oscar Daniel Flores Linares}
\newcommand{\ayudanteTeoriaDos}{Claudia Itzel Gutiérrez Sánchez}

\newcommand{\ayudanteLabUno}{Ricardo Badillo Macías}
\newcommand{\ayudanteLabDos}{Rocío Aylin Huerta González}

\newcommand{\equipoNombre}{Dobles Comillas}

\newcommand{\integranteUno}{Contreras Morales Román Ismael}
\newcommand{\integranteDos}{Cornejo de la mora Iñaki}
\newcommand{\integranteTres}{Hipólito Felipe Christopher Guillermo}
\newcommand{\integranteCuatro}{Ledesma Cuevas Eduardo}

% ==========================================

\begin{document}

\thispagestyle{empty}

\begin{center}
	\vspace*{1cm}

	\begin{tabular}{cc}
		
\includegraphics[width=3cm]{escudos/unam_seal.pdf} & 
\includegraphics[width=3cm]{escudos/fc_seal.pdf}
	\end{tabular}

	\vspace{0.8cm}
	{\LARGE\bfseries UNIVERSIDAD NACIONAL AUTÓNOMA DE MÉXICO}\\[0.3cm]
	{\large FACULTAD DE CIENCIAS, 2027-I}\\[0.3cm]
	{\large\bfseries FUNDAMENTOS DE BASES DE DATOS}

	\vspace{1.2cm}
	{\huge\bfseries \documentoTitulo:}\\[0.3cm]
	{\Large \documentoSubtitulo}

	\vspace{1.5cm}
\end{center}

\begin{multicols}{2}
	\raggedcolumns

	\textbf{PROFESOR:}\\
	\profesorNombre\\[1.2cm]

	\textbf{AYUDANTES DE TEORÍA:}\\
	\ayudanteTeoriaUno\\
	\ayudanteTeoriaDos\\[1.2cm]

	\textbf{AYUDANTES DE LABORATORIO:}\\
	\ayudanteLabUno\\
	\ayudanteLabDos\\[1.2cm]

	\columnbreak

	\textbf{EQUIPO:}\\
	\equipoNombre\\[1.2cm]

	\textbf{INTEGRANTES DEL EQUIPO:}\\
	\integranteUno\\
	\integranteDos\\
	\integranteTres\\
	\integranteCuatro

\end{multicols}

\vspace{1cm}

% ==========================================
% Modelo Entidad-Relación Extendido
% ==========================================
\section*{Modelo Entidad-Relación Extendido}

% >>> PROPUESTA (a revisar)
% Introducción y justificación de alcance. No proviene de Explicación.pdf.
El presente documento describe y justifica el modelo entidad-relación extendido (EER)
construido a partir de las especificaciones del caso de uso del centro de entretenimiento
familiar ``PuellaGame''. Para el modelado se emplearon los conceptos del modelo EER
---entidades, atributos, relaciones, cardinalidades, participación y generalización/
especialización--- tal como se presentan en \parencite{silberschatz2006}. Cada decisión de
diseño se justifica con base en los requerimientos del caso, indicando explícitamente
cuando una justificación proviene del documento previo \emph{Explicación.pdf}
\marcaOriginal{} y cuando se trata de una mejora propuesta por el equipo
\marcaPropuesta{}.
% <<< PROPUESTA (a revisar)

% ==========================================
% Diagrama
% ==========================================
\subsection*{Diagrama Entidad-Relación}

\begin{figure}[htbp]
	\makebox[\textwidth][c]{\includegraphics[width=1.1\textwidth]{Diagramas/diagramaPractica02.drawio.png}}
	\caption{Modelo entidad-relación extendido del centro de entretenimiento familiar PuellaGame.}
	\label{fig:diagrama_er}
\end{figure}

% ==========================================
% Justificación de las relaciones (Explicación.pdf)
% ==========================================
% >>> FUENTE: Explicación.pdf
\subsection*{Justificación de las relaciones \marcaOriginal}

A continuación se justifica cada relación del diagrama indicando su cardinalidad y la
participación de cada extremo.

\begin{description}[leftmargin=0pt, style=nextline]

	\item[ADMINISTRAR --- Gerente $\leftrightarrow$ Sucursal \textnormal{(1:1)}]
		Un gerente administra a lo sumo una sucursal y una sucursal es administrada por un
		único gerente. La participación del gerente es \textbf{total}: todo gerente debe
		tener una sucursal asignada, pues no hay gerente sin sucursal. La participación de
		la sucursal es \textbf{parcial}, ya que una sucursal puede existir antes de
		asignarle un gerente.

	\item[PERTENECE --- Máquina $\leftrightarrow$ Sucursal \textnormal{(N:1)}]
		Por especificaciones del caso de uso, una máquina pertenece a una única sucursal,
		pero una sucursal tiene muchas máquinas. La participación de la máquina es
		\textbf{total}, porque toda máquina registrada en el inventario debe estar
		localizada en una sucursal. La participación de la sucursal es \textbf{parcial},
		pues una sucursal puede darse de alta sin tener máquinas instaladas todavía.

	\item[CONTENER --- Máquina $\leftrightarrow$ Juego \textnormal{(N:1)}]
		Un mismo juego puede instalarse en muchas máquinas distintas, pero cada máquina de
		arcade ejecuta un único juego dedicado. La participación de la máquina es
		\textbf{total}, pues no hay máquinas vacías: cada máquina debe tener un juego para
		operar. La participación del juego es \textbf{parcial}, pues puede existir un juego
		en el catálogo que actualmente no esté instalado en ninguna máquina física.

	\item[TENER --- Cliente $\leftrightarrow$ Tarjeta \textnormal{(1:1)}]
		Por especificaciones del caso de uso, a cada cliente se le asocia una sola tarjeta
		activa. La participación del cliente es \textbf{total}: un cliente no existe si no
		se le ha emitido su tarjeta, pues para obtenerla se debe registrar al cliente con
		sus datos. La participación de la tarjeta también es \textbf{total}, pues toda
		tarjeta emitida debe estar vinculada a los datos de un cliente.

	\item[EMITIR --- Cajero $\leftrightarrow$ Tarjeta \textnormal{(1:N)}]
		Un solo cajero puede emitir muchas tarjetas. La participación del cajero es
		\textbf{parcial}, pues un cajero puede tener cero tarjetas emitidas. La
		participación de la tarjeta es \textbf{total}, pues toda tarjeta requiere el
		registro del cajero que la dio de alta.

	\item[EFECTUAR --- Sucursal $\leftrightarrow$ Recarga \textnormal{(1:N)}]
		En una sucursal se pueden realizar muchas recargas. La participación de la recarga
		es \textbf{total}, pues una recarga no puede realizarse sin una sucursal. La
		participación de la sucursal es \textbf{parcial}, pues una sucursal puede existir
		sin un registro de recarga.

	\item[RECIBIR --- Tarjeta $\leftrightarrow$ Recarga \textnormal{(1:N)}]
		Una tarjeta puede recibir varias recargas. La participación de la recarga es
		\textbf{total}, pues se necesita una tarjeta para poder realizar la recarga. La
		participación de la tarjeta es \textbf{parcial}, pues podemos tener una tarjeta sin
		haber realizado ninguna recarga.

	\item[PROCESAR --- Cajero $\leftrightarrow$ Recarga \textnormal{(1:N)}]
		Un cajero puede procesar varias recargas. La participación de la recarga es
		\textbf{total}, pues para procesar una recarga se necesita un cajero que la
		realice. La participación del cajero es \textbf{parcial}, pues puede existir un
		cajero sin haber procesado una recarga.

	\item[OCURRIR --- Máquina $\leftrightarrow$ Partida \textnormal{(1:N)}]
		En una máquina ocurren muchas partidas. La participación de la partida es
		\textbf{total}, pues para que ocurra una partida se necesita una máquina. La
		participación de la máquina es \textbf{parcial}, pues una máquina puede existir sin
		haber sido jugada aún.

	\item[REGISTRAR --- Tarjeta $\leftrightarrow$ Partida \textnormal{(1:N)}]
		Una tarjeta puede ser utilizada para muchas partidas. La participación de la
		partida es \textbf{total}, pues para registrar una partida es necesario registrar
		con qué tarjeta se pagó. La participación de la tarjeta es \textbf{parcial}, pues
		puede existir una tarjeta sin registrar ninguna partida aún.

	\item[SOLICITAR --- Cliente $\leftrightarrow$ Canje \textnormal{(1:N)}]
		Un cliente puede solicitar varios canjes. La participación del canje es
		\textbf{total}, pues se necesita un cliente que solicite cualquier canje. La
		participación del cliente es \textbf{parcial}, pues un cliente puede no haber
		solicitado canjes.

	\item[INCLUIR --- Premio $\leftrightarrow$ Canje \textnormal{(1:N)}]
		Un tipo de premio puede canjearse varias veces. La participación del canje es
		\textbf{total}, pues en cada canje hay un registro de qué premio se canjeó. La
		participación del premio es \textbf{parcial}, pues un premio puede aún no ser
		canjeado.

	\item[PROCESAR --- EncargadoDePremios $\leftrightarrow$ Canje \textnormal{(1:N)}]
		El encargado de premios puede procesar varios canjes. La participación del canje es
		\textbf{total}, pues en cada canje es necesario registrar al empleado que lo
		procesó. La participación del empleado de premios es \textbf{parcial}, pues puede
		existir sin haber procesado un canje.

	\item[OCURRIR --- Sucursal $\leftrightarrow$ Canje \textnormal{(1:N)}]
		En una sucursal pueden ocurrir varios canjes. La participación del canje es
		\textbf{total}, pues en cada canje hay que tener un registro de la sucursal donde
		ocurrió. La participación de la sucursal es \textbf{parcial}, pues puede existir
		sin haber ocurrido un canje.

	\item[REALIZAR --- Técnico $\leftrightarrow$ Mantenimiento \textnormal{(1:N)}]
		Un técnico puede realizar muchos mantenimientos. La participación del
		mantenimiento es \textbf{total}, pues en cada mantenimiento debe haber un registro
		del técnico que lo realizó. La participación del técnico es \textbf{parcial}, pues
		un técnico puede no haber hecho ningún mantenimiento aún.

	\item[RECIBIR --- Máquina $\leftrightarrow$ Mantenimiento \textnormal{(1:N)}]
		Una máquina puede recibir muchos mantenimientos. La participación del
		mantenimiento es \textbf{total}, pues en cada mantenimiento se debe registrar qué
		máquina lo recibió. La participación de la máquina es \textbf{parcial}, pues una
		máquina puede existir sin ningún mantenimiento aún.

	\item[REGISTRAR --- Sucursal $\leftrightarrow$ Inventario \textnormal{(1:N)}]
		Una sucursal puede tener varias filas en la tabla de inventario. La participación
		del inventario es \textbf{total}, pues es necesario saber la sucursal de donde
		proviene cada producto. La participación de la sucursal es \textbf{parcial}, pues
		una sucursal puede no tener productos disponibles.

	\item[CONTENER --- Premio $\leftrightarrow$ Inventario \textnormal{(1:N)}]
		Un mismo premio puede estar en varias filas del inventario, pues puede estar
		disponible en distintas sucursales. La participación del inventario es
		\textbf{total}, pues es necesario saber qué premio está disponible. La
		participación del premio es \textbf{parcial}, pues pueden existir premios dados de
		alta pero sin estar disponibles aún en ninguna sucursal.

\end{description}

% ==========================================
% Generalización y especialización (Explicación.pdf)
% ==========================================
\subsection*{Generalización y especialización \marcaOriginal}

El caso de uso requiere modelar distintos roles operativos para el personal. Para evitar la
redundancia de atributos compartidos y mantener la consistencia, se utilizó una abstracción
de generalización:

\begin{itemize}
	\item \textbf{Superentidad:} \emph{Empleado}. Concentra los atributos comunes (RFC,
	      CURP, nombre, contacto, salario y fecha de contratación).
	\item \textbf{Subentidades:} \emph{Cajero}, \emph{Gerente}, \emph{Técnico} y
	      \emph{Encargado de Premios}.
	\item \textbf{Restricción de completitud (total):} toda entidad registrada como empleado
	      debe pertenecer forzosamente a al menos uno de los roles operativos específicos;
	      no existen empleados genéricos sin rol.
	\item \textbf{Restricción sobre el carácter disjunto (disjunta):} un empleado no puede
	      pertenecer a más de una subentidad al mismo tiempo, pues cada rol está bien
	      definido y separado.
\end{itemize}
% <<< FUENTE: Explicación.pdf

% ==========================================
% Mejoras propuestas
% ==========================================
\begin{cajapropuesta}

	\textbf{\large Mejoras propuestas \marcaPropuesta}

	Los apartados siguientes son aportaciones del equipo sobre \emph{Explicación.pdf} y pueden
	eliminarse de forma independiente sin afectar el resto del documento.

	\subsection*{Justificación de las entidades \marcaPropuesta}

	Cada entidad del diagrama responde a un requerimiento explícito del caso de uso:

	\begin{itemize}
		\item \emph{Sucursal}, \emph{Máquina} y \emph{Juego}: el caso pide almacenar las
		      sucursales, las máquinas (con su estado) y el catálogo de juegos que posee el
		      centro.
		\item \emph{Cliente} y \emph{Tarjeta}: se requiere registrar al cliente para asociarle
		      una tarjeta de acceso (una sola activa) con la que paga y acumula puntos.
		\item \emph{Recarga} y \emph{Partida}: se debe llevar registro de cada recarga y de cada
		      partida (puntaje, puntos ganados, fecha y hora).
		\item \emph{Premio}, \emph{Inventario} y \emph{Canje}: se deben administrar los premios,
		      su disponibilidad por sucursal y cada canje efectuado.
		\item \emph{Mantenimiento}: se debe registrar cada mantenimiento con fecha, tipo y
		      descripción del problema.
		\item \emph{Empleado} (con sus subentidades): el caso exige registrar al personal y sus
		      roles (gerente, cajero, técnico y encargado de premios).
	\end{itemize}

	\subsection*{Justificación de los atributos \marcaPropuesta}

	\begin{itemize}
		\item \textbf{Compuestos:} \emph{NombreCompleto} (nombres, apellido paterno, apellido
		      materno) y \emph{Dirección} (calle, número interior, número exterior, colonia,
		      estado) agrupan datos que siempre se consultan juntos.
		\item \textbf{Multivaluados:} \emph{Correos}, \emph{Teléfonos} y
		      \emph{Certificaciones}, ya que el caso permite varios valores por registro.
		\item \textbf{Derivados:} \emph{Edad} se obtiene de \emph{FechaNacimiento} y
		      \emph{esVIP} se deduce del conteo de visitas del jugador; \emph{CantidadDisponible}
		      se deriva de los movimientos de inventario.
		\item \textbf{Claves:} \emph{RFC} y \emph{CURP} identifican de forma natural a un
		      empleado; el resto de las entidades usa identificadores sustitutos
		      (\emph{idCliente}, \emph{idTarjeta}, \emph{idPartida}, \emph{idCanje},
		      \emph{idRecarga}, \emph{idMantenimiento}) que evitan ambigüedad y facilitan las
		      referencias.
	\end{itemize}

	\subsection*{Entidades débiles y asociativas \marcaPropuesta}

	\emph{Inventario} no posee clave propia: queda identificado por la pareja
	(\emph{Sucursal}, \emph{Premio}) y su único atributo descriptivo es
	\emph{CantidadDisponible}, por lo que puede modelarse como entidad débil dependiente de
	\emph{Sucursal}. Las entidades \emph{Canje}, \emph{Partida}, \emph{Recarga} y
	\emph{Mantenimiento} cuentan con identificador sustituto propio y funcionan como entidades
	asociativas que materializan el histórico de operaciones.

	\subsection*{Reglas del negocio reflejadas en el modelo \marcaPropuesta}

	\begin{itemize}
		\item Cada jugador puede tener únicamente una tarjeta activa (relación TENER 1:1).
		\item La emisión de la tarjeta cuesta \$20.00 MXN y exige una recarga inicial de
		      \$10.00 MXN; las recargas y los costos de juego son múltiplos de 10.
		\item Los puntos dependen del tipo de juego: 1 punto por acierto en Destreza y Disparos,
		      y 20 puntos por quedar en los primeros 5 lugares en Velocidad y Ritmo.
		\item La categoría VIP se otorga con 10 visitas en un mismo mes, da 10\,\% de puntos
		      adicionales y se conserva repitiendo el mínimo de visitas.
		\item El canje requiere un mínimo de 20 puntos y el premio debe corresponder al rango de
		      edad del jugador (Infantil 3--12, Juvenil 13--17, Adulto 18+).
		\item Los premios se clasifican por puntos: Bajos (20--1{,}000), Medios
		      (1{,}001--3{,}999) y Grandes (4{,}000--10{,}000+).
	\end{itemize}

\end{cajapropuesta}

\printbibliography

\end{document}
